\chapter{The Learner's World Already Has a Physics}
\label{sec:causality-intro}
% ===================================================================

The learner of the preceding part is not a mind contemplating a world.
It is a term sitting in parallel with the terms it wants to know about,
made of the same material, subject to the same rules, paying the same
prices.
That was the point of building it that way.
But it has a consequence which has so far gone unremarked, and this part
is about the consequence.

If the learner is a term in a cost-accounted theory, then it is already
inside something that deserves to be called a physics.
It has states, and laws that take one state to another, and a currency it
must spend to move.
Nobody put a physics there.
The physics is what the theory looks like from inside, and every
cost-accounted interactive theory has one.

The question this part asks is therefore not ``how do we get a physics
out of computation'' --- that was the question the earlier arrangement of
this material asked, and it presupposed a knower standing outside,
building.
The question is: \emph{which} physics, and why that one?
A learner opening its eyes in a $\catciGSLT$ finds itself in a world with
some structure and not other structure, and the structure it finds
depends on which axioms the theory happens to satisfy.
Most objects of $\catciGSLT$ satisfy very few of them.
The world humanity actually turns out to live in satisfies a
surprisingly long list, and each item on that list cuts the category
down to a subcategory.

\section{Physics as a sequence of subcategories}
\label{sec:physics-ladder}

The organizing device of this part is a ladder.
At the top sits $\catciGSLT$, containing every cost-accounted interactive
theory whatever.
Each chapter below imposes one further axiom, and each axiom names a full
subcategory whose objects are the theories satisfying it.
A learner in a theory low on the ladder inhabits a world that looks
increasingly like the one physicists describe; a learner high on the
ladder inhabits a world in which most of what physics says is simply
false --- not because the physics is wrong, but because its hypotheses do
not hold there.

\begin{center}
\renewcommand{\arraystretch}{1.3}
\begin{tabular}{>{\itshape}p{4.4cm} p{5.4cm} c}
\toprule
\normalfont\textbf{Axiom imposed} & \textbf{What the learner then has} & \textbf{Ch.} \\
\midrule
none
  & states, laws, a budget
  & --- \\
every rewrite has an inverse
  & a time-symmetric world, with the arrow of time in the boundary
    condition rather than the laws
  & \ref{sec:reversibility} \\
histories are recorded
  & a causal order, and a proper time
  & \ref{sec:causal} \\
rewrites carry a semiring value
  & a decoration --- the weighting of Chapter~\ref{ch:weight} with its
    codomain chosen: nondeterminism, cost, probability or amplitude,
    according to the semiring
  & \ref{sec:weighted} \\
conversion is arbitrage-free
  & a single conserved quantity, and a first law
  & \ref{sec:cost} \\
the decoration is internalised
  & dynamics as a fact \emph{in} the logic, not merely about it
  & \ref{sec:extended-modal} \\
values are complex
  & interference, and an obstruction
  & \ref{sec:quantum} \\
\bottomrule
\end{tabular}
\end{center}

Two things should be said about this table before it is unpacked.

The first is that the axioms are genuinely independent in the sense that
matters: there are theories satisfying any given one and failing the
next, and Chapter~\ref{sec:synthesis} exhibits some.
The ladder is not a chain of definitions dressed up as discoveries.

The second is that the ladder does not descend all the way to
\emph{our} physics, and this part does not claim that it does.
It descends some distance and then stops at an obstruction which
Chapter~\ref{sec:quantum} states precisely: the complex instance gives
amplitudes and gives interference, and it does not by itself give the
Born rule.
What is offered is a family of possible physics with our own somewhere
inside it, not a derivation of ours.

\begin{remark}[Two ladders, and they are perpendicular]
\label{rmk:two-ladders}
Chapter~\ref{ch:tower}, in the Prestige, describes a different ladder ---
a tower of oracle extensions, in which each position extends the ones
below it and the physics available at a higher position is strictly
richer.
It is easy to confuse the two, and they are perpendicular.
The ladder of this part moves \emph{down} within a fixed computational
power, adding axioms and thereby shrinking the class of theories under
consideration.
The ladder of Chapter~\ref{ch:tower} moves \emph{up}, adding computational
power and thereby enlarging what any given theory can express.
A theory can be low on this ladder and low on that one --- a very
classical world inhabited by very limited agents --- or high on both.
The two coordinates are independent, and several confusions about what
computation can and cannot ground come from collapsing them into one.
\end{remark}

\section{What the earlier arrangement got wrong}

It is worth recording, since some readers will have seen an earlier form
of this material.

That presentation carried two maps rather than one, called them a
Lagrangian and a Hamiltonian, and claimed that everything else followed
by construction.
The first was an unnecessary duplication: one decoration into a semiring
suffices, and the choice of semiring does the work the two maps were
introduced to do.
The second was a misnaming.
The third was an overstatement, and the most interesting of the three
errors, because the parts that do \emph{not} follow by construction turn
out to be exactly the parts that carry content.
Conservation holds only under a condition; that condition has a
cohomological measure; and the identification of the conserved charge
with the generator of time evolution remains a conjecture.
The table below is annotated accordingly, and the annotations should be
taken seriously.

\section{The Rosetta Stone}
\label{sec:rosetta}

The following previews the correspondence developed in this part.
The third column records the \emph{status} of each row:
\textsc{con} marks a construction, true by definition of the objects
involved; \textsc{thm} marks a theorem with a hypothesis that can fail;
\textsc{cnj} marks a conjecture we have not proved.
A fully annotated version appears in Chapter~\ref{sec:synthesis}.

\begin{center}
\renewcommand{\arraystretch}{1.35}
\begin{tabular}{>{\itshape}p{4.3cm} p{5.9cm} c}
\toprule
\normalfont\textbf{GSLT} & \textbf{Physics} & \textbf{Status} \\
\midrule
Term                          & State / initial condition & \textsc{con} \\
Rewrite rule                  & Law of motion & \textsc{con} \\
Rewrite path                  & Trajectory & \textsc{con} \\
Synchronization tree          & Causal graph & \textsc{con} \\
Minimum path length           & Proper time & \textsc{con} \\
Equations $\eqs$              & Gauge equivalence & \textsc{con} \\
Reversible envelope $\GSLT^\dagger$ & Time-symmetric theory & \textsc{con} \\
Decoration, tropical          & Action functional; least action & \textsc{thm} \\
Decoration, complex           & Path amplitude $e^{iS/\hbar}$ & \textsc{con} \\
Sum over paths                & Feynman path integral & \textsc{con} \\
Interference of paths         & Quantum interference & \textsc{cnj} \\
No-arbitrage conversion       & Existence of an energy function & \textsc{thm} \\
Virtual token $\nu$           & Energy, as a Noether charge & \textsc{thm} \\
Holonomy of the energy form   & Failure of energy to be a state function & \textsc{thm} \\
Dissipation $\theta$          & Second law; free energy & \textsc{thm} \\
Ledger law $\sigma+\kappa=\sigma_0$ & Potential plus kinetic & \textsc{thm} \\
Erasure bounded by holonomy   & Landauer's principle & \textsc{thm} \\
Located authority             & Entanglement of space and time & \textsc{cnj} \\
$\nu$ generates reduction     & Hamiltonian, generator of evolution & \textsc{cnj} \\
Symmetry of the decoration    & Noether current & \textsc{cnj} \\
\bottomrule
\end{tabular}
\end{center}

The reader who has come through Part~\ref{part:ecology} will notice that
several rows of this table have already been met wearing other clothes.
The no-arbitrage row is Theorem~\ref{thm:vtok-ftap}, introduced in
Part~\ref{part:machinery} as a fact about exchange rates and about to be
read as a fact about energy.
The ledger law is the conservation identity the mortal scientist lives
under: what it has not yet spent, plus what it has spent, is what it
started with.
The second law is why a scientist that is merely right can still starve.
This is neither coincidence nor analogy.
It is the same accounting, and the reason the learner had to be built
first is that the accounting is easier to believe once one has watched
something live under it.

\section{Organization of this part}

Chapter~\ref{sec:reversibility} constructs the time-symmetric envelope
$\GSLT^\dagger$ and locates the arrow of time in the boundary condition.
Chapter~\ref{sec:causal} identifies the synchronization trees with causal
graphs in the sense of Bombelli, Lee and Sorkin, and reads a proper time
off the minimum path length.
Chapter~\ref{sec:weighted} takes the weighting endofunctor of
Chapter~\ref{ch:weight} --- one map into a semiring --- and asks what the
choice of semiring delivers, locating the least-action principle in the
tropical instance.
Chapter~\ref{sec:cost} asks what is conserved and at what price.
It is the longest chapter in the part, and it contains the first law, the
second law, the ledger law and the Landauer bound; it also re-derives the
virtual token of Chapter~\ref{ch:vtok} in its physical reading, which is
the one place in the book where a single theorem is deliberately proved
twice.
Chapter~\ref{sec:extended-modal} extends the modal operators to carry the
full dynamical state of a transition, so that the dynamics becomes a fact
stated in the logic rather than merely about it.
Chapter~\ref{sec:quantum} lifts the decoration to complex amplitudes,
states the obstruction honestly --- there are two of them, and only one is
about bisimulation --- and finds that a presentation has a second place to
put coherence, namely the equations it declines to quotient.
Chapter~\ref{sec:synthesis} collects the ladder into one picture, and
Chapter~\ref{sec:discussion} says what this part does not establish.

% ===================================================================
